%======================================================================
\section{미리 정한 규칙과 그 기록}
\label{sec:protocol}

\begin{yoji}
규칙 A와 결합 연산자는 이 연구가 로그를 하나도 추출하기 전에, 규칙 B는 W1 로그를 질의하기 8초 전에 커밋했어요. 규칙 B의 가설은 W0의 탐색적 비교에서 나왔어요. 수정 분석(달러 금액, 다섯 토큰, 동결 직전 12개월 점수 기간)은 등록 결과를 모두 안 뒤에 정했으므로, 그 판정은 재분석이에요.
\end{yoji}

표~\ref{tab:timeline}는 약속을 한 순서대로, 각 약속의 기록과 그때 이미 알던 것을 보여 줘요. 해석에 중요한 점은 세 가지예요.
\begin{enumerate}
\item 규칙 A와 결합 연산자는 이 연구의 로그를 하나도 추출하기 전에 정했어요. 그때는 이 연구의 점수도 라벨도 하나도 계산되지 않았어요.
\item 규칙 B의 동기가 된 비교는 W0의 라벨을 안 뒤에 그 창에서 눈에 띄었어요. 그래서 그 창의 증거는 탐색적이에요.
\item 수정 분석(달러 금액, 다섯 토큰, 동결마다 직전 1년의 점수 기간)은 미리 정한 분석의 결과를 모두 안 뒤에 정했어요. 같은 규칙과 기준을 적용하지만, 그 판정은 등록한 시험이 아니에요. \ref{sec:results}절은 수정 분석을 먼저, 등록 분석을 그 옆에 보고해요.
\end{enumerate}

\begin{table}[!ht]
\centering
\caption{약속의 순서, 각 약속의 기록, 약속할 때 이미 알던 것(2026년, 시각은 UTC)}
\label{tab:timeline}
\footnotesize
\setlength{\tabcolsep}{3pt}
\renewcommand{\arraystretch}{1.1}
\begin{tabularx}{\linewidth}{@{}L{2.2cm}YL{2.5cm}L{3.3cm}@{}}
\toprule
언제 & 약속이나 사건 & 기록 & 이미 알던 것 \\
\midrule
10월 9일 13:41:27 & 분석 계획과 설정을 커밋: 코호트 규칙, 점수(C-PR, $\lambda\in\{0.25,0.5,0.75\}$와 그 끝점, S-PR), W0 라벨, 통계, 규칙~A; 금액은 원래 단위, 토큰은 모든 ERC-20 & \sha{43d8045} & 이 연구의 로그를 하나도 추출하지 않음 \\
10월 9일 13:41:37--13:43:54 & 관측 기간의 스캔~A; BigQuery 작업은 계획 커밋 10초 뒤에 만들어짐 & 작업 메타데이터, 비용 장부 & --- \\
10월 9일 16:49 전 & 계약 27{,}844개의 코호트, 에고 추출, 같은 기간과 W0 평가; 규칙~A 판정 & 요약을 \sha{0d3a057}에 커밋 & 관측 기간의 로그 \\
10월 9일 16:49:32 & W1 등록: $\delta=0.01$의 규칙~B, F4--F6, 민감도 분석 S1과 S2, 짝 연구의 트레이더 결과 비교(그 규칙 R1--R3과 결론 문구); 설정의 상태를 registered로 바꿈 & \sha{0d3a057} & 등록 분석의 같은 기간과 W0 결과 모두; 트레이더의 W0 이익 라벨(탐색적 실행) \\
10월 9일 16:49:33 & 추출 스크립트가 커밋된 등록을 확인하고 매니페스트를 씀(커밋 1초 뒤): 커밋과 SHA-256 해시 두 개 & 매니페스트 & --- \\
10월 9일 16:49:40--16:50:11 & 2026년 7--9월의 스캔~B; 작업은 커밋 8초 뒤에 만들어짐; 매니페스트는 16:50:55에 닫힘 & 작업 메타데이터, 매니페스트 & 그 전에 이 연구를 위해 추출한 7--9월 로그 없음 \\
10월 9일 18:00:55 & 등록 W1 평가, 이어서 강건성 점검과 시빌 모형을 커밋 & \sha{8d8b568} & W1 라벨 \\
10월 9일 23:24:39 & 등록 벤치마크와 자료 설명(BigQuery 작업 시각, 금액 비율)을 커밋 & \sha{f33a87a} & W1 라벨 \\
나중 & 수정 분석, 계획과 등록에서 벗어난 변경: 다섯 토큰(WETH, USDC, USDT0, ARB, WBTC)의 달러 금액, 기울기 $b=\ln 3/m$, C-PR 허락 층의 상한 $U$, 동결 직전 12개월의 점수와 동결마다 코호트 하나; 상수 $m_T=170$, $m_A=100$, $U=19{,}000$달러; 규칙~A, B, R1--R3을 그 기준대로 적용; 결과 커밋 & 커밋은 표~\ref{tab:ids}, 결과는 \sha{950a7a4} & 등록 결과 모두 \\
나중 & 부록~\ref{app:explore}의 탐색적 점검 & 이 원고의 보조 자료 & 위의 결과 모두 \\
\bottomrule
\end{tabularx}

\vspace{4pt}
\parbox{\linewidth}{\footnotesize 커밋 시각은 커밋한 사람의 시각이에요. UTC+9로 기록된 것을 UTC로 바꿨어요. 작업 시각은 BigQuery가 작업마다 알려 주는 생성 시각과 끝 시각이에요. 커밋의 전체 식별자와 파일 해시는 표~\ref{tab:ids}에 있어요.}
\end{table}

\subsection{규칙 A(분석 계획, 2026년 10월 9일 13:41:27 UTC)}

규칙 A는 C-PR이 두 층을 각각 혼자 걷는 것보다 나은지를 물어요.
\begin{itemize}
\item \textbf{주 부분}은 첫 홀드아웃 창 W0(점수는 2026년 3월 31일에 동결, 라벨은 2026년 4–6월)을 써요. $\lambda=0.5$에서 C-PR은 두 가지를 해내야 해요.
  \begin{itemize}
  \item A1: 새 허락 쌍에서 허락 걷기보다 더 나쁘지 않아야 해요. 곧 $\Delta\tau=\tau_{\text{C-PR}}-\tau_{\lambda=1}$의 95\% 구간이 0을 포함하거나 0보다 위에 있어야 해요.
  \item A2: 새 송금자에서 송금 걷기를 이겨야 해요. 곧 $\tau_{\text{C-PR}}-\tau_{\lambda=0}$의 구간이 0보다 위에 있어야 해요.
  \end{itemize}
\item \textbf{부 부분}은 점수 자체의 기간(같은 기간)에서 매칭 코호트를 써요. C-PR의 점 추정값이 송금 점검 가족과 허락 점검 가족에서 각각 더 나은 한 층의 95\% 구간 안에 들어야 해요. 그리고 시빌 보정 안정성 점검 세 가지의 가족에서는 두 층을 모두 넘어야 해요. 시빌 보정 안정성은 계획이 붙인 이름이고, 들어오는 상대 가운데 서로 다른 상대의 몫, 활동 기간, 활동한 달 수로 이뤄져요.
\end{itemize}
주 부분이 실패하면 “C-PR이 한 층만 걷는 걷기를 이기는가”의 답은 ‘아니요’예요. 그리고 라벨을 본 뒤에 $\lambda$를 조정하지 않기로 했어요.

규칙의 두 기준은 설계부터 약해요.
\begin{itemize}
\item 정의대로의 “더 나쁘지 않음”은 정밀하지 못할수록 유리해요. 구간이 넓을수록 0을 포함하기 쉽기 때문이에요.
\item 부 부분은 짝지은 대비를 쓰지 않고, 한 점수의 점 추정값을 다른 점수의 주변 구간과 견줘요.
\end{itemize}
이 절차를 다시 쓰는 사람은 앞의 것을 비열등성 허용 폭으로, 뒤의 것을 짝지은 대비로 바꿔야 해요. 이 원고는 규칙 옆에 짝지은 대비를 함께 보고해요.

\begin{puri}[판정 기준 세 가지]
\begin{itemize}
\item \textbf{우월성:} 차이 $\Delta\tau$의 95\% 구간 전체가 0보다 위에 있어야 해요. 예를 들어 F1의 +0.170 [0.149, 0.189]는 충족이에요.
\item \textbf{비열등성:} 구간의 아래 끝이 $-\delta=-0.01$보다 위에 있어야 해요. 조금 뒤지는 것은 봐주되, 0.01보다 더 뒤질 수 있으면 안 된다는 뜻이에요. F2의 −0.025 [−0.038, −0.015]는 구간 전체가 −0.01보다 아래라 실패예요.
\item \textbf{더 나쁘지 않음:} 구간이 0을 포함하거나 0보다 위면 충족이에요. A1의 +0.031 [0.021, 0.042]는 충족이에요. 이 기준은 구간이 넓기만 해도 통과하기 쉬워요.
\end{itemize}
\end{puri}

\subsection{규칙 B(등록, 2026년 10월 9일 16:49:32 UTC)}
\label{sec:rule-b}

규칙 B는 두 번째 창 W1에서 C-PR이 송금 걷기보다 나은지를 $\Delta\tau=\tau_{\text{C-PR}}-\tau_{\lambda=0}$로 물어요. W1은 2026년 6월 30일에 동결하고, 라벨은 2026년 7월 1일부터 9월 30일까지 세요.
\begin{itemize}
\item F1(스펜더, 새 허락 쌍), 우월성: $\Delta\tau$의 95\% 구간 전체가 양수여야 해요.
\item F2(스펜더, 새 송금자), 비열등성: 구간의 아래 끝이 $-\delta$보다 위에 있어야 해요.
\item F3(트레이더, 청산 없이 닫은 비율), 같은 허용 폭의 비열등성.
\end{itemize}
셋이 모두 성립해야 해요. 판정이 세 조건의 교집합이므로, 각각을 5\% 수준에서 시험해도 개선을 잘못 주장할 확률은 다중 비교 보정 없이 5\% 이하로 유지돼요.

허용 폭 $\delta=0.01$은 F2가 되풀이하는 대비의 W0 구간에서 왔어요. 그 대비는 −0.030, [−0.034, −0.026]이고 반폭이 0.0040이며, 이것을 소수 둘째 자리로 올림했어요. 그러니 이 값은 스펜더 대비의 정밀도를 잴 뿐, 실제로 중요하다고 판단한 차이가 아니에요. 등록은 이 허용 폭을 트레이더 대비 F3에도 적용했어요. 그런데 F3의 구간은 훨씬 적은 계약에 기대요(\ref{sec:precision}절). 등록은 F3의 약점도 미리 적어 두었어요. 어느 걷기도 청산을 맞히지 못해 두 계수가 모두 0 근처여도 F3은 통과할 수 있어요. 그래서 판정 옆에 두 계수를 함께 보고해요.

\begin{puri}[왜 셋을 다 요구하면 보정이 필요 없나요?]
개선했다는 주장이 틀렸다면, 세 조건 가운데 적어도 하나는 실제로 성립하지 않아요. 그런데 판정이 통과하려면 그 조건의 시험도 통과해야 해요. 그 시험은 5\% 수준이라, 성립하지 않는 조건을 통과시킬 확률이 5\%를 넘지 않아요. 그래서 개선을 잘못 주장할 확률도 5\%를 넘지 않아요.
\end{puri}

판정 밖에서, 정해진 순서로 보고할 결과도 세 가지 정했어요.
\begin{itemize}
\item F4: 청산 없이 닫은 비율에서 C-PR이 송금 걷기보다 우월한지.
\item F5: 같은 라벨에서 허락 걷기보다 우월한지.
\item F6: 규칙 A의 주 부분을 W1에서 다시 하기.
\end{itemize}
“정해진 순서”는 보고하는 순서일 뿐이에요. 앞 결과가 통과해야 다음을 보는 관문 순서(gatekeeping)는 없어요. 그래서 F4–F6은 보정하지 않은 95\% 구간으로 본 명목상의 결과예요. 민감도 분석 두 개도 함께 등록했어요.
\begin{itemize}
\item S1: AWP를 비교 대상으로 해서 F1–F4를 다시 해요.
\item S2: 그래프가 빠뜨린 코호트 계약을 0점 대신 외톨이 점으로 두고 F1–F6을 다시 해요.
\end{itemize}
둘 다 판정에 들어가지 않아요. 점수, 풀이 프로그램, 질의, 코호트 필터, 거래소 사건 해독기, 부트스트랩(짝지은 다시 뽑기 400번, 시드 42)은 계획의 것을 그대로 썼어요.

\subsection{추출의 잠금}
\label{sec:gate}

설정 파일에는 W1의 상태 칸이 있고, 등록 커밋이 이 칸을 “registered”로 바꿨어요. 추출 스크립트는 무엇을 질의하기 전에 두 가지를 확인해요. 등록 문서와 설정이 로컬 변경 없이 커밋되어 있는지, 그리고 상태가 “registered”인지예요. 그러고 나서 커밋과 두 파일의 SHA-256 해시를 매니페스트에 적어요.
\begin{itemize}
\item 잠금은 해시를 비교하지 않아요. 해시는 독자가 나중에, 커밋된 파일이 스크립트가 본 파일과 같은지 확인하게 해 줄 뿐이에요.
\item 스크립트에는 우회 장치가 없어요. 매니페스트의 \texttt{override\_used} 칸은 아무것도 기록하지 않는 상수예요.
\item 표~\ref{tab:ids}은 커밋의 전체 식별자와 해시를 보여 줘요. 매니페스트는 저자들의 Windows 작업본에서 CRLF 줄 끝으로 파일을 해시했어요. git이 저장하는 대로 LF 줄 끝인 같은 글은 둘째 해시 쌍을 주고, 두 쌍 모두 커밋된 파일과 대조했어요.
\end{itemize}

기록이 무엇을 보여 주는지는 누가 그 기록을 쥐고 있는지에 달려 있어요.
\begin{itemize}
\item 커밋 시각은 커밋한 기계의 시계가 정하고, 커밋한 사람이 고를 수 있어요.
\item BigQuery 작업 시각(스캔~A의 생성 13:41:37, 스캔~B의 생성 16:49:40, 표~\ref{tab:timeline})은 클라우드 제공자가 프로젝트의 작업 기록에 보관하고, 프로젝트 구성원만 읽을 수 있어요. 저자들은 이 시각을 \sha{f33a87a}(10월 9일 23:24:39 UTC)에 커밋한 자료 설명에 옮겨 적었어요.
\end{itemize}
그래서 바깥사람은 등록의 내용, 해시, 커밋의 순서를 확인할 수 있어요. 하지만 등록이 W1 추출보다 먼저였다는 주장은 커밋 시각과 작업 시각에 기대고, 제3자의 시간 도장에 기대지 않아요. 추출 전에 커밋 해시를 체인 위나 독립된 등록처에 박아 두는 것 같은 믿을 만한 시간 도장이 있으면 이 빈틈이 사라져요.

\begin{puri}[커밋으로 잠근다는 것은 무엇인가요?]
등록 파일을 커밋하지 않았거나 커밋한 뒤 작업본에서 고쳐 두었으면, W1 로그를 추출하는 스크립트는 아예 돌지 않아요. 그러니 스크립트가 W1 로그를 뽑았다면 그때 등록 파일은 커밋된 상태였다는 뜻이에요. 다만 그 커밋이 \emph{언제}였는지는 커밋한 사람의 시계와 BigQuery의 작업 기록에 기대요.
\end{puri}

\begin{table}[!ht]
\centering
\caption{미리 정한 기록과 수정 분석의 커밋 식별자(시각은 UTC), 그리고 등록 파일과 계획의 SHA-256 해시}
\label{tab:ids}
\footnotesize
\setlength{\tabcolsep}{0pt}
\begin{tabular}{@{}>{\raggedright\arraybackslash}p{0.8\linewidth}@{}}
\toprule
계획과 설정(규칙 A), 10월 9일 13:41:27\\
\quad\sha{43d80457142dc002c4410d4af1c3c164f39a3272}\\
W1 등록(규칙 B, R1--R3), 10월 9일 16:49:32\\
\quad\sha{0d3a057e4f0c8473d36d89fbe33c24deff0c9911}\\
등록 결과 커밋, 10월 9일 18:00:55\\
\quad\sha{8d8b568ecaf882cccc6caef60b444823966f5ce8}\\
등록 벤치마크와 자료 설명, 10월 9일 23:24:39\\
\quad\sha{f33a87a92f3f894c829451e55cc7ff9a33d74e3f}\\
\midrule
\emph{수정 분석(달러 금액, 다섯 토큰, 동결 직전 12개월의 점수), 등록 결과를 안 뒤에 정함}\\
달러 금액, 토큰 규칙, 기울기\\
\quad\sha{a8974759f8a8e6e0c8b07c23481ff429eaa492ee}\\
다섯 토큰과 가격 입력\\
\quad\sha{09dfc323de2c3f0ffa40c7a6a0c5adf0b94c1122}\\
C-PR 허락 층의 상한 규칙\\
\quad\sha{2705b28a4653d47a6940cca064bf923562db8f9d}\\
점수 기간: 동결 직전 12개월\\
\quad\sha{e6c352dcbc5ba14e7bcd25df7bdfeefedb1091ce}\\
코호트와 상수, 결과의 코드 리비전\\
\quad\sha{0a9421563a529762662ccaef4d86ac96bfad03f6}\\
수정 결과 커밋\\
\quad\sha{950a7a4cd9401002b3be9348797b4a8dd9d9f075}\\
\midrule
등록 문서의 SHA-256, CRLF 줄 끝(매니페스트에 기록)\\
\quad\sha{2126c8dc3fe3c0f9273aa4ecfc1e15bd}\\
\quad\sha{dc1f5a6f65e98548b9528f2b18ca2a6b}\\
등록 문서의 SHA-256, LF 줄 끝(git이 저장하는 대로)\\
\quad\sha{34cd6dfeb56b8fb8a553cb24afbff06c}\\
\quad\sha{6a94e6f69a2eb7f7d642b0cc39e54b2c}\\
설정의 SHA-256, CRLF 줄 끝(매니페스트에 기록)\\
\quad\sha{d6a30adbcabcf511d6892332da73639e}\\
\quad\sha{fdbece25891b392e13e90faf09b7a833}\\
설정의 SHA-256, LF 줄 끝(git이 저장하는 대로)\\
\quad\sha{2ea235021241fa408f8776d55c2b594d}\\
\quad\sha{34ced0ca1dad44e16f22775c5db8297f}\\
\sha{43d8045}에 있는 계획의 SHA-256, LF 줄 끝\\
\quad\sha{7ec80d49c17331bd10e086fecd909a3f}\\
\quad\sha{36020dbb284f0dbfbb81a8549f84fd9f}\\
\bottomrule
\end{tabular}

\vspace{4pt}
\parbox{0.8\linewidth}{\footnotesize SHA-256 값은 하나를 두 줄로 나눠 적었어요. 띄어 쓰지 않고 이어 붙이면 돼요.}
\end{table}

\subsection{수정 분석}
\label{sec:revision}

모든 창의 등록 결과를 안 뒤, 계획이 무게를 다는 방식에서 결함 세 가지가 드러났어요.
\begin{enumerate}
\item 원래 단위 위에서 $b=1$이면, 송금의 99.52\%와 마지막 허락의 99.78\%에서 $V(z)$와 1의 차이가 $10^{-4}$보다 작았어요. 그래서 금액이 엔도스랭크와 AWP에서 아무 역할을 하지 못했어요.
\item C-PR의 각 층은 원래 단위를 토큰 사이에서 그냥 더했어요. 그런데 USDC 1개는 원래 단위로 $10^{6}$개, 감싼 이더 1개는 $10^{18}$개예요. 그리고 무제한 허락($2^{256}-1$ 원래 단위, 마지막 허락의 45.3\%)이 그 허락을 가진 주인의 행을 거의 다 차지했어요.
\item 모든 ERC-20 계약의 로그가 그래프에 들어왔어요. 그런데 토큰 계약은 어떤 주소든 주인이나 보낸 쪽으로 적은 로그를, 그 주소가 실제로 움직였든 아니든 쓸 수 있어요.
\end{enumerate}

수정 분석은 이 결과를 보면서, 달러 금액, 다섯 토큰, 동결 직전 12개월 점수 기간으로 정의했어요. 자세히는 이래요.
\begin{itemize}
\item 모든 금액을 DefiLlama 가격 계열의 하루 종가로 달러 값으로 매겨요~\cite{defillamanddocs}.
\item 토큰은 다섯 개를 남겨요. 위험도가 낮음이나 보통인 Chainlink 시장 가격 피드로 가격이 매겨지는 토큰 가운데~\cite{chainlinknddatafeeds}, 2026년 6월 30일 전 12개월에 행이 가장 많은 다섯 개예요(WETH, USDC, USDT0, ARB, WBTC).
\item 모든 점수를 동결 직전 12개월의 사건으로 만들고, 동결마다 코호트를 하나씩 둬요. 동결마다 1년을 쓰므로 1년보다 오래된 사건은 점수에도, 같은 기간 점검에도, 라벨을 정할 때 보는 이력에도 들어가지 않아요. 그리고 두 동결을 같은 길이의 기간으로 점수 매겨요.
\item 엔도스랭크와 AWP에는 $b=\ln 3/m$을 써요. 중앙값은 W0 동결까지 12개월에서 구해서, 그 뒤의 사건이 들어가지 않게 해요.
\item C-PR 허락 층의 허락마다 상한 $U=19{,}000$달러를 둬요. 같은 달들의 다섯 토큰 단일 송금의 99번째 백분위(18,575달러)를 유효숫자 두 자리로 반올림한 값이에요. 유한 허락의 백분위는 상한으로 쓸 수 없어요. 그 99번째 백분위가 $10^{63}$달러 근처라, 무제한 허락이 여전히 그 주인들의 행을 차지하기 때문이에요.
\end{itemize}
모든 코호트, 그래프, 대리 지표, 스펜더 라벨을 이미 추출한 로그에서 다섯 토큰과 12개월 기간으로 계산해요. 규칙 A와 B는 그 기준, $\delta$, $\lambda$를 그대로 적용해요. 달러 금액, 다섯 토큰, 12개월 점수 기간을 쓰는 수정 분석은 등록 결과를 보면서 정한, 계획과 등록에서 벗어난 변경이에요. 표~\ref{tab:ids}이 그것을 정의하는 커밋을 보여 주고, 그 점수에 내린 판정은 모두 재분석이에요.

\begin{puri}[왜 원래 단위로는 금액이 아무 역할을 못 했나요?]
원래 단위는 소수점 조정을 하지 않은 가장 작은 단위예요. USDC 1개는 원래 단위로 $10^{6}$개, 감싼 이더 1개는 $10^{18}$개예요. $b=1$이면 원래 단위로 10개만 넘어도 $V(z)$와 1의 차이가 $10^{-4}$보다 작아져서, 거의 모든 송금과 허락이 같은 무게를 받았어요. 반대로 금액을 그대로 더하는 C-PR의 층에서는 숫자가 터무니없이 큰 무제한 허락이 그 주인의 행을 거의 혼자 차지했어요.
\end{puri}

\subsection{미리 정하지 않은 것}

A1, A2, F1–F6, S1, S2 밖의 대비는 모두 그 창의 라벨을 안 뒤에 계산했어요.
\begin{itemize}
\item W0에서 새 허락 쌍으로 C-PR을 송금 걷기와 견준 것.
\item S-PR의 대비들.
\item 원시 차수와의 대비 모두.
\item S1 밖에서 C-PR을 엔도스랭크나 AWP와 견준 것 모두.
\end{itemize}
트레이더의 W0 이익 라벨은 탐색적 실행에서 계산했고, 그 계수는 등록과 함께 커밋했어요. W0의 청산 라벨은 이 연구의 어떤 점수에 대해서도 계산하지 않았어요.

계획 커밋에는 계획, 설정, 로그를 복사하는 질의, BigQuery 실행기가 들어 있어요. 코호트, 에고 네트워크, 쌍 표, W0 라벨을 만드는 질의와 스크립트, 모든 점수와 부트스트랩의 구현, 같은 기간과 W0 평가를 돌리는 스크립트는 등록 커밋과 함께 저장소에 들어갔어요. 관측 기간에서 이미 돌고 규칙 A를 판정한 뒤였으므로, 그 전의 코드는 기록에 없어요. \ref{sec:model}절의 단위 시험도 그때 함께 들어갔어요. W1 추출과 평가는 커밋된 그대로의 것을 썼어요. 부록~\ref{app:explore}의 점검은 모두 이 모든 일이 끝난 뒤에 돌렸어요.

%======================================================================
\section{자료와 측정}
\label{sec:data}

\begin{yoji}
로그는 아비트럼 원의 공개 BigQuery 표에서 왔어요. 수정 분석은 동결마다 직전 12개월의 다섯 토큰 사건으로 코호트와 그래프를 만들어요. 수정 분석의 매칭 코호트는 6,483개, 스펜더 코호트는 W0 7,914개와 W1 7,960개이고, W1 트레이더는 두 분석 모두 70개예요.
\end{yoji}

\subsection{출처와 코호트}

모든 로그는 이더리움 위의 낙관적 롤업인 아비트럼 원(체인 ID 42161)에서 왔어요~\cite{kalodner2018arbitrum,thibault2022rollups}. 읽은 곳은 공개 BigQuery 표예요.
\begin{center}
\texttt{\small bigquery-public-data.goog\_blockchain\_arbitrum\_one\_us.logs}
\end{center}
\begin{itemize}
\item 스캔~A는 관측 기간인 2023년 10월 1일부터 2026년 6월 30일까지의 로그를, 스캔~B는 2026년 7–9월의 로그를 복사했어요.
\item 수정 분석에서는 점수마다 그 동결에서 끝나는 1년을 써요. W1과 같은 기간 점검은 2025년 7월부터 2026년 6월까지, W0은 2025년 4월부터 2026년 3월까지예요.
\item 등록 분석은 관측 기간 안에서 동결까지의 모든 사건으로 점수를 매겨요.
\item 나이는 동결에서 거꾸로 재요.
\item 허락과 송금 로그는 ERC-20 사건 서명과 정확히 세 개의 토픽으로 골라요. 그러면 ERC-721 토큰이 빠져요. 금액이 0인 송금은 버리고, 금액이 0인 허락은 남겨요. 허락을 지우는 사건이기 때문이에요.
\item 닫기 하나는 GMX V2 사건 발신 계약이 남긴 \texttt{PositionDecrease}라는 이름의 \texttt{EventLog1} 기록이에요. 일부 닫기와 청산을 똑같이 세요. 주문 종류가 청산이거나, 청산 처리 계약이 남긴 기록이면 청산으로 세요~\cite{gmxnddocs}.
\end{itemize}

매칭 코호트에는 점수 기간에 서로 다른 주인 셋 이상에게서 0이 아닌 허락을 받은 계약이 모두 들어가요.
\begin{itemize}
\item 수정 분석에서는 동결마다 직전 12개월의 다섯 토큰 허락으로 세요. 2026년 6월 30일 동결에서 6,483개(매칭 코호트), W0 동결에서 6,343개이고, 그 가운데 5,383개가 둘 다에 있어요.
\item 계획의 코호트는 관측 기간의 모든 토큰 허락으로 세요. 이 규칙을 만족한 스펜더 주소 29,022개 가운데, 블록 513,216,506에서(\texttt{eth\_getCode}로 읽음) 27,844개가 계약 코드를 가졌고, 1,121개는 코드가 없었고, 57개는 EIP-7702 위임 표시를 가졌어요~\cite{buterin2024eip7702}. 뒤의 두 무리는 외부 소유 계정이라 뺐어요.
\item 수정 코호트 둘은 모두 계획의 코호트 안에 있어요.
\end{itemize}
한쪽 끝에 코호트 계약이 있는 사건만 추출해요. 허락은 코호트 계약이 주인이나 스펜더인 것, 송금은 코호트 계약이 보낸 쪽이나 받는 쪽인 것이에요. 그래서 각 그래프는 코호트의 한 걸음 에고 네트워크이고, 수정 분석에서는 그 동결의 코호트의 에고 네트워크예요. 2026년 6월 30일 전 12개월 동안 수정 에고 표에는 허락 로그 85,986,562개와 송금 로그 442,722,039개가 있어요. 같은 달들에 계획의 에고 표가 가진 로그의 63.3\%와 76.2\%예요. 관측 기간 전체로는 계획의 에고 표에 389,866,882개와 1,538,404,067개가 있어요.

\begin{puri}[에고 네트워크란?]
코호트 계약에 한쪽 끝이 닿는 사건만 모은 그래프예요. 코호트 밖의 주소끼리 오간 허락과 송금은 들어오지 않아요. 그래서 나가는 무게의 합, 다시 출발의 무게, 걷는 사람이 갈 수 있는 길은 체인 전체가 아니라 이 에고 네트워크의 것이에요.
\end{puri}

\subsection{그래프}
\label{sec:graphs}

2026년 6월 30일 동결 시점의 두 그래프는 이래요.
\begin{itemize}
\item 허락 그래프: 수정 분석에서 점 2,514,846개와 화살표 3,808,949개, 등록 분석에서 점 13,055,150개와 화살표 22,321,470개.
\item 송금 그래프: 수정 분석에서 점 3,603,300개와 화살표 8,426,124개, 등록 분석에서 점 23,029,298개와 화살표 59,737,617개.
\end{itemize}
수정 분석의 점 수에는, 먼지 같은 금액만 오간 쌍의 끝점 일부가 들어 있어요. 그런 쌍은 상한이 있는 무게 $V$가 부동소수점에서 0이 돼요. 화살표 수는 그런 쌍을 빼고 세고, 풀이 프로그램도 빼요. 그래서 엔도스랭크의 그래프는 점이 2,513,982개, AWP의 그래프는 3,601,727개예요. 반면 C-PR과 S-PR의 금액 무게 층은 그런 쌍을 남기므로, 그 송금 층의 화살표는 8,432,635개예요(표~\ref{tab:bench}). 등록 분석에서는 0이 되는 무게가 없어서 수가 같아요.

양수인 마지막 허락 가운데 무제한 허락은 수정 분석에서 49.8\%이고, 상한 $U$에 닿는 것은 55.7\%예요. 등록 분석에서는 무제한 허락이 45.3\%예요. 어떤 걷기의 그래프에 없는 코호트 계약은 그 걷기에서 0점이에요(S2가 이것을 바꿔요). 수정 분석에서 허락 그래프 밖에 191개, 송금 그래프 밖에 1,192개, 둘 다의 밖에 30개가 있어요.

\subsection{코호트와 라벨}

창마다의 스펜더 코호트에는, 점수 기간의 허락 자료에서 동결 시점에 양수의 마지막 허락을 가진 계약이 모두 들어가요. 코호트 계약이 허락한, 매칭 코호트 밖의 계약도 들어가요. 라벨은 두 가지예요.
\begin{itemize}
\item \textbf{새 허락 쌍:} 라벨 기간에 그 스펜더에게 양수의 허락을 준 주인 가운데, 동결 시점에 그 쌍의 허락이 양수가 아니었던 주인의 수예요.
\item \textbf{새 송금자:} 스펜더 말고 라벨 기간에 그 스펜더에게 송금을 보낸 주소 가운데, 점수 기간의 시작부터 동결까지는 보낸 적이 없는 주소의 수예요.
\end{itemize}
W0의 스펜더 코호트는 수정 분석에서 7,914개, 등록 분석에서 33,101개이고, W1에서는 7,960개와 34,587개예요. 표~\ref{tab:holdout}가 각 라벨을 얻은 스펜더의 수를 보여 줘요.

W1 트레이더 코호트는 GMX V2의 계약 계정 가운데, 어느 쪽이든 동결일 그래프에 나오고 2026년 7–9월에 포지션을 셋 이상 닫은 것이에요. 두 분석 모두 70개이고, 그 가운데 9개는 청산이 한 번 이상 있었고, 1개는 닫기가 모두 청산이었어요. 라벨은 청산 없이 닫은 비율, 곧 1에서 라벨 기간의 닫기 가운데 청산으로 끝난 몫을 뺀 값이에요. 어느 층으로도 만들지 않았고, 비율이라 거래를 많이 한다고 유리해지지 않아요. 계약 계정은 관측 기간에 포지션 133,685개를 닫았고(그 가운데 청산 6,473개), 2026년 3분기에는 6,805개를 닫았어요(청산 182개).

\subsection{통계}

모든 계수는 동점을 바로잡은 켄달의 $\tau_b$이고~\cite{kendall1938rank}, SciPy로 계산해요. 구간은 코호트의 계약을 다시 뽑는 짝지은 부트스트랩으로 구한 95\% 백분위 구간이에요~\cite{efron1979bootstrap,diciccio1996bootstrap}.
\begin{itemize}
\item 다시 뽑기의 번호는 NumPy의 기본 생성기에서 시드 42로 뽑아요. 코호트마다 번호 행렬을 하나 두고, 한 라벨의 모든 점수에 같은 행렬을 써요. 그래서 두 점수의 차이는 모두 짝지은 차이예요.
\item 모든 결과는 다시 뽑기 400번을 써요.
\item 보조 스크립트는 S2를 뺀 모든 판정 대비를 발표된 계약별 점수 파일에서 다시 계산하고, 발표된 값을 정확히 재현해요(부록~\ref{app:explore}).
\end{itemize}
